\lyseis{7.1}{Θετικοί και αρνητικοί αριθμοί}
\sol{1} α) $-5$ \quad β) $+12$ \quad γ) $-8$ \quad δ) $+15$ \quad ε) $-2$ \quad στ) $-30$.
\sol{2} Θετικοί: $+3{,}5$, $8$, $+0{,}01$. Αρνητικοί: $-7$, $-\frac{1}{2}$. Το $0$ δεν είναι ούτε θετικός ούτε αρνητικός.
\sol{3} α) Ομόσημοι. β) Ετερόσημοι. γ) Ομόσημοι. δ) Ετερόσημοι.
\sol{4} Ακέραιοι: $-4$, $0$, $\frac{9}{3}=3$, $10$.
\sol{5} Ο $-4$ είναι $4$ cm αριστερά του $0$, ο $-1{,}5$ είναι $1{,}5$ cm αριστερά, ο $2$ $2$ cm δεξιά και ο $3{,}5$ $3{,}5$ cm δεξιά.
\sol{6} $A(-4)$, $B(-1{,}5)$, $\Gamma(1)$, $\Delta(3{,}5)$.
\sol{7} α) Λάθος. β) Σωστό. γ) Σωστό. δ) Λάθος: ομόσημοι. ε) Σωστό. στ) Σωστό.
\sol{8} Αθήνα $+7$, Θεσσαλονίκη $-3$, Ιωάννινα $-5$, Χανιά $+11$, Φλώρινα $-10$. Στον κατακόρυφο άξονα, από πάνω προς τα κάτω: $+11$, $+7$, $0$, $-3$, $-5$, $-10$.
\sol{9} Από το $3$ στο $-2$ κατεβαίνει $5$ ορόφους ($3,\,2,\,1,\,0,\,-1,\,-2$).
\sol{10} Το ψάρι ($-45$ μ.) είναι πιο κοντά στην επιφάνεια ($0$). Το αεροπλάνο είναι πάνω από την επιφάνεια, στα $300$ μ.
\sol{11} Μυρτώ $+20$, Στέλιος $-7$, Φανή $+5$, Βασίλης $-12$.
\sol{12} Από το $-3$ ανεβαίνω $5$ μονάδες δεξιά: $-2,\,-1,\,0,\,1,\,2$. Η θερμοκρασία γίνεται $+2$°.
\sol{13} Παράδειγμα: ένας χειμωνιάτικος θερμοκρασία $-4$°C, ή ο λογαριασμός τράπεζας με $-50$ €, ή το βάθος της θάλασσας.
\sol{14} Το $0$ είναι το σημείο που χωρίζει τις δύο πλευρές: δεν είναι ούτε στη θετική ούτε στην αρνητική. Δεν έχει κατεύθυνση.
\sol{15} Το μηδέν είναι σύμβαση: ανάλογα με την κλίμακα, το $0$°C αντιστοιχεί στα $32$°F. Δεν σημαίνει «καθόλου θερμοκρασία».
\sol{16} Με τα άκρα: $-5,\dots,5$, δηλαδή $11$. Χωρίς τα άκρα $9$.
\sol{17} Ελεύθερη απάντηση. Π.χ. θερμοκρασία, υψόμετρο και βάθος, όροφοι, λογαριασμός τράπεζας, έτη π.Χ.
\sol{18} Ελεύθερη απάντηση.
